\clearpage
\begin{explanationbox}
\textbf{先看一次具体翻折：}
\begin{center}
\begin{tikzpicture}[x=1cm,y=1cm]
\draw[p055 frame] (0,0) rectangle (13.8,7.5);
\node[anchor=west,font=\small\bfseries] at (.4,7) {竖直镜面：左右反号，上下不变};
\begin{scope}[shift={(6.9,5.45)},scale=1.1]
  \draw[p055 mirror] (0,-3.65)--(0,.6);
  \node[teal!70!black,right] at (.15,-2.1) {$x=2$};
  \coordinate (P) at (0,0); \coordinate (A) at (-2,-3); \coordinate (B) at (2,-3);
  \draw[p055 original,-Stealth] (P)--(A);
  \draw[p055 image,-Stealth] (P)--(B);
  \draw[CStatement,-Stealth] (P)--(-2,0) node[midway,above] {$\Delta x=-2$};
  \draw[CStatement,-Stealth] (-2,0)--(A) node[midway,left] {$\Delta y=-3$};
  \draw[orange!85!black,-Stealth] (P)--(2,0) node[midway,above] {$\Delta x'=2$};
  \draw[orange!85!black,-Stealth] (2,0)--(B) node[midway,right] {$\Delta y'=-3$};
  \node[above] at (0,.55) {$P(2,1)$};
  \fill (P) circle (1.7pt);
  \fill[CStatement] (A) circle (1.7pt) node[below left] {$A(0,-2)$};
  \fill[orange!85!black] (B) circle (1.7pt) node[below right] {$A'(4,-2)$};
\end{scope}
\node at (6.9,.75) {$k=\dfrac{-3}{-2}=\dfrac32\qquad\longrightarrow\qquad k'=\dfrac{-3}{2}=-\dfrac32$};
\end{tikzpicture}

\end{center}
从 $P$ 到 $A$，横、纵增量都是负数，商却为正；沿 $x=2$ 翻折后，只有横增量变号，商就反号。图中的两根箭头等长，这是轴对称保持长度的结果。

\medskip
\textbf{推广到任意方向：} 把图中的 $(-2,-3)$ 换成 $(p,q)$。竖直镜面给出 $(-p,q)$，水平镜面给出 $(p,-q)$，所以都得到 $-q/p$。$45^\circ$ 镜面把横纵增量互换（可能同时反号），所以得到 $p/q$。

\medskip
\textbf{为什么求斜率可以少算一步？} 斜率只看两个坐标差的比值。镜面的平移不改变这个比值；同向放大或缩小整根箭头，也不改变比值。但这只是比较\emph{方向}的简化，实际反射仍保持原箭头长度。

\medskip
\textbf{任意斜镜面：} 将原方向绕镜面翻到另一侧，与镜面的夹角大小保持相等。方法三先画等角图，再从方向角关系推导斜率；不必死记通式。

\medskip
\textbf{内心题的识别信号：} 顶点 $P$ 与内心 $I$ 同在 $x=a$ 上，说明这条竖直线平分顶角的\emph{内角}，因此两边射线镜像，有限斜率互为相反数。例1会画出这个构型，并解释为什么两条边不必等长。

\medskip
\textbf{最后区分两个问题：} “像直线朝哪个方向”看增量；“像直线到底在哪里”还要找经过的点。射线、位置和竖直边界的具体对照分别见例题与易错图。
\end{explanationbox}

\clearpage
\begin{explanationbox}
\textbf{再从角度理解（一）：竖直角平分线}
\begin{center}
% P=(0,3), A=(-2,0), B=(2,0); theta=atan(2/3).
\begin{tikzpicture}[x=1cm,y=1cm]
\draw[p055 frame] (0,0) rectangle (13.8,7.6);
\node[anchor=west,font=\small\bfseries] at (.4,7.1) {竖直对称：两个倾斜角关于 $\frac\pi2$ 对称};
\begin{scope}[shift={(5.2,2)},scale=1.15]
  \coordinate (P) at (0,3); \coordinate (A) at (-2,0); \coordinate (B) at (2,0);
  \draw[p055 axis] (-3.4,0)--(3.6,0) node[right] {$x$};
  \draw[p055 mirror] (0,-.5)--(0,3.8);
  \node[teal!70!black,anchor=west] at (.2,3.65) {$x=0$};
  \draw[p055 original] (A)--(P);
  \draw[p055 image] (P)--(B);
  \draw[gray] (.18,0)--(.18,.18)--(0,.18);
  \fill (P) circle (1.8pt) node[above left] {$P$};
  \fill[CStatement] (A) circle (1.8pt) node[below] {$A$};
  \fill[orange!85!black] (B) circle (1.8pt) node[below] {$B$};
  \node[below left] at (0,0) {$H$};
  \draw (P) ++(236.309932:.7) arc (236.309932:270:.7);
  \draw (P) ++(270:.85) arc (270:303.690068:.85);
  \node at (-.36,1.95) {$\theta$};
  \node at (.4,1.95) {$\theta$};
  \draw[CStatement,-Stealth] (A) ++(0:.7) arc (0:56.309932:.7);
  \node[CStatement] at (-1.05,.48) {$\alpha$};
  \draw[orange!85!black,-Stealth] (B) ++(0:.8) arc (0:123.690068:.8);
  \node[orange!85!black] at (2.55,1.05) {$\beta$};
\end{scope}
\node[anchor=west,align=left] at (9.65,4.9) {$0<\theta<\frac\pi2$\\[9pt]$\alpha=\frac\pi2-\theta$\\[9pt]$\beta=\frac\pi2+\theta$};
\node at (6.9,.65) {从水平正向逆时针量倾斜角：$\alpha+\beta=\pi$。};
\end{tikzpicture}

\end{center}
图取 $P(0,3),A(-2,0),B(2,0)$，$PH$ 平分 $\angle APB$，两侧小角均为 $\theta$。直线的\textbf{倾斜角}是从水平正向逆时针量到直线向上方向的角，范围为 $[0,\pi)$；斜率等于倾斜角的正切。

\medskip
由直角三角形，$PA$ 的倾斜角为 $\alpha=\frac\pi2-\theta$；$PB$ 在 $B$ 处向左上方延伸，因此倾斜角为 $\beta=\frac\pi2+\theta=\pi-\alpha$。于是
\[
k_{PA}=\tan\!\left(\frac\pi2-\theta\right),\qquad
k_{PB}=\tan(\pi-\alpha)=-\tan\alpha=-k_{PA}.
\]
\textbf{关键观察：} 顶角的两个 $\theta$ 相等，转成倾斜角却是一锐一钝，不能把它们都当成 $\theta$。图中 $k_{PA}=3/2,\ k_{PB}=-3/2$；平移镜面，或沿两边射线另取点，都不改变这一方向关系。
\end{explanationbox}

\clearpage
\begin{explanationbox}
\textbf{再从角度理解（二）：水平角平分线}
\begin{center}
% P=(0,0), A=(4,2), B=(4,-2); theta=atan(1/2).
\begin{tikzpicture}[x=1cm,y=1cm]
\draw[p055 frame] (0,0) rectangle (13.8,7.6);
\node[anchor=west,font=\small\bfseries] at (.4,7.1) {水平对称：下方射线的锐角不是直线的倾斜角};
\begin{scope}[shift={(3.2,3.55)}]
  \coordinate (P) at (0,0); \coordinate (A) at (4,2); \coordinate (B) at (4,-2);
  \draw[p055 axis] (0,-2.5)--(0,2.7) node[above] {$y$};
  \draw[p055 mirror,-Stealth] (-1.5,0)--(5,0) node[right] {$x$};
  \draw[p055 helper] (A)--(B);
  \draw[p055 original] (P)--(A);
  \draw[p055 image] (P)--(B);
  \fill (P) circle (1.8pt) node[below left=4pt] {$P=O$};
  \fill[CStatement] (A) circle (1.8pt) node[above right] {$A(s,t)$};
  \fill[orange!85!black] (B) circle (1.8pt) node[below right] {$B(s,-t)$};
  \draw (0:.85) arc (0:26.565051:.85);
  \draw (0:1.05) arc (0:-26.565051:1.05);
  \node at (1.35,.35) {$\theta$};
  \node at (1.5,-.35) {$\theta$};
  \draw[p055 axis] (B)--(5.8,-2);
  \draw[orange!85!black,-Stealth] (B) ++(0:.75) arc (0:153.434949:.75);
  \node[orange!85!black] at (4.35,-.95) {$\beta$};
  \node[teal!70!black,anchor=east] at (-.2,.35) {$y=0$};
\end{scope}
\node[anchor=west,align=left] at (9.8,4.5) {$s,t>0$\\[9pt]$\tan\theta=t/s$\\[9pt]$\beta=\pi-\theta$};
\node at (6.9,.6) {$PA$ 的倾斜角为 $\theta$，$PB$ 的倾斜角为 $\pi-\theta$。};
\end{tikzpicture}

\end{center}
图取 $P=O$、$A(s,t)$、$B(s,-t)$，其中 $s,t>0$。两条射线与水平轴的锐角均为 $\theta$，但 $PB$ 向右下方的方向对应角 $-\theta$，它不在倾斜角的取值范围内。

\medskip
沿直线 $PB$ 取向左上方的方向，倾斜角就是 $\beta=\pi-\theta$。图中在 $B$ 处补画水平正向，直接标出这个钝角，故
\[
k_{PA}=\tan\theta=\frac ts,\qquad
k_{PB}=\tan(\pi-\theta)=-\tan\theta=-\frac ts.
\]
\textbf{关键观察：} 相等的是两侧锐角；两条直线的倾斜角互补，所以正切反号。也可写 $k_{PB}=\tan(-\theta)$，但不能把 $-\theta$ 称为倾斜角。此图取 $s=4,t=2$；平移到任意水平镜面 $y=b$ 后，斜率关系仍相同。
\end{explanationbox}
